
\begin{enumerate}
  \item 半径 $1$ の円に内接する $6$ 個の半径の等しい円を図 $1$ のように描き，さらに図 $2$ のように $6$ 個の小さな半径の等しい円を描く，この操作を無限にくり返したとき，$6$ 個ずつ次々に描かれる円の面積の総和 $S_2$ と，それらの円の円周の長さの総和 $C_2$ を求めよ．
  \item (1) で $6$ 個の円を次々に描いていった．一般に，自然数 $n \ge 2$ に対して $3n$ 個の円を用いて同様の操作を行うとき，描かれる円の面積の総和 $S_n$ と，それらの円の円周の長さの総和 $C_n$ を求めよ．
  \item 数列 $S_2, S_3, S_4, \cdots$ の極限値を求めよ．
\end{enumerate}

\begin{figure}[H]
  \centering
  \begin{subcaptionblock}{0.5\linewidth}
    \centering
    \begin{tikzpicture}
      % 中心に配置する円（オプション: 問題文には明示されていないが、図1には中心に円がない）
      % \draw (0,0) circle (0.5cm);

      % 周囲の6つの円
      % 外側の円の半径をR=3cmとする。
      % 周囲の円の半径をrとする。
      % 中心から周囲の円の中心までの距離をdとする。
      % d = R - r
      % また、中心から周囲の円の中心までの距離がd、周囲の円の半径がrであるとき、
      % sin(30度) = r / d より、d = 2r
      % よって、R - r = 2r => R = 3r => r = R/3
      % d = 2R/3
      \def\R{3} % 外側の円の半径
      \tikzmath{
        \r = \R / 3; % 周囲の円の半径
        \d = \R - \r; % 中心から周囲の円の中心までの距離
      }
      % 外側の大きな円
      \draw (0,0) circle (\R); % 適当なサイズに調整
      \foreach \i in {0, 60, ..., 300} {
          \draw ({\d*cos(\i)}, {\d*sin(\i)}) circle (\r cm);
        }
    \end{tikzpicture}
    \subcaption{図1}
  \end{subcaptionblock}\hfill
  \begin{subcaptionblock}{0.5\linewidth}
    \centering
    \begin{tikzpicture}
      \def\R{3} % 外側の円の半径
      \tikzmath{
        \r = \R / 3; % 周囲の円の半径
        \d = \R - \r; % 中心から周囲の円の中心までの距離
      }
      % 最も外側の大きな円（半径1と仮定）
      \draw (0,0) circle (\R);
      % 外側の大きな円
      \foreach \i in {0, 60, ..., 300} {
          \coordinate (C\i) at ({\d*cos(\i)}, {\d*sin(\i)});
          \draw (C\i) circle (\r);
          \draw (0,0) -- (C\i); % 中心から外側の円の中心への直線
        }
      \tikzmath{
        \rsmall = \r / 3; % 周囲の円の半径
        \dsmall = \r - \rsmall; % 中心から周囲の円の中心までの距離
      }
      % 内側の小きな円
      \foreach \i in {0, 60, ..., 300} {
          \draw ({\dsmall*cos(\i)}, {\dsmall*sin(\i)}) circle (\rsmall);
        }
    \end{tikzpicture}
    \subcaption{図2}
  \end{subcaptionblock}\hfill
\end{figure}